\chapter{1957:Lord Todd}

\label{chap:chemistry1957}

The Nobel Prize in Chemistry for the year 1957 was awarded to Lord Todd.

\textbf{Motivation:}

for his work on nucleotides and nucleotide co-enzymes

\section{Lord Todd}

\subsection*{Early Life and Education}

Lord Todd was born on 1907-10-02 in Glasgow, Scotland.

\subsection*{Career and Major Research}

Alexander Robertus Todd studied at the University of Glasgow, took his doctorate at the University of Frankfurt, and earned a doctorate at the University of Oxford. He held the chair of chemistry at the University of Manchester and in 1944 became professor of organic chemistry at the University of Cambridge, where he remained for the rest of his career. He worked out the structures of the nucleotides and of the nucleoside coenzymes, and his group achieved the first syntheses of adenosine triphosphate (ATP), adenosine diphosphate (ADP), and the coenzyme flavin-adenine dinucleotide (FAD).

\subsection*{Nobel Prize-winning Work}

The work for which Lord Todd was awarded the Nobel Prize in Chemistry in 1957 was described by the Nobel Committee as: "for his work on nucleotides and nucleotide co-enzymes".

Todd established how nucleotides are linked through their sugar and phosphate units and showed that the coenzymes are built from nucleotides joined in specific arrangements. He also contributed to the chemistry of vitamin B12 and of the nucleic acids, predicting a linear phosphate-sugar backbone for RNA.

\subsection*{Significance and Impact}

Todd's nucleotide chemistry provided the structural basis for understanding ATP as the energy currency of cells and for the later understanding of nucleic acid structure. His syntheses of the coenzymes proved their structures beyond doubt and made them available for biochemical study.

\subsection*{Later Career and Honors}

Todd was raised to the peerage as Lord Todd of Trumpington in 1962 and served as master of Christ's College, Cambridge, and president of the Royal Society of London. He died on 1997-01-10 in Cambridge, England.

\subsection*{Legacy}

Lord Todd is regarded as the founder of modern nucleotide chemistry, and his work laid the structural foundations for the discovery of the double helix and for the chemistry of nucleic acids.

\subsection*{Selected Publications}

\begin{itemize}

\item \emph{Synthesis of Adenosine Triphosphate}, Nature, 1948.

\end{itemize}

\clearpage

\section*{Chapter Summary}

The 1957 prize recognized work on nucleotides and nucleotide coenzymes. Lord Todd determined the structures of the nucleotides and coenzymes and achieved their chemical synthesis at Cambridge, providing the structural basis for understanding ATP and the nucleic acids.

